\section{Quotation marks and Science-bar}
\label{sec:quotation}

\begin{orientacao}{Desenvolvimento posterior --- sete lugares reservados}
Esta parte recebe a pergunta deixada pelo guardanapo. Os blocos abaixo
são espaços de planejamento, não novos atos já redigidos. A estrutura
preserva sua proposta de sete atos por parte. Substitua cada bloco por
um ambiente \texttt{ato} e escreva o trecho acadêmico correspondente.
\end{orientacao}

\atoplanejado{I}{Retomar a circulação do guardanapo. Desenvolver o
estrangeirismo enunciativo e a mudança de registro em diálogo com
\refobra{cunha2017}. Especificar quem toma emprestado o nome e quais
obrigações a tradução produz.}

\atoplanejado{II}{Alterar uma regra do bar para interrogar a pretensão
de uma Ciência que encerra a discussão. Apresentar Science-bar como
proposta do ensaio e precisar sua relação com a quarta garantia da
Constituição moderna: \refobra{latour1994}, pp.~38--39. Explicar
também os limites dessa analogia.}

\atoplanejado{III}{Dar corpo e localização a Cecinepá. Explicitar se é
um lugar existente, uma experiência, um dispositivo de escrita ou
uma combinação desses modos. Desenvolver sua condição de heterotopia
com \refobra{foucault2013}, sem fazer de todo espaço alternativo uma
heterotopia automaticamente.}

\atoplanejado{IV}{Fazer aparecer uma situação em que coisas, pessoas,
palavras e obrigações não possam ser inventariadas isoladamente.
Definir a proposta de \textit{total ontological fact} e distinguir
seu alcance do fato social total de Mauss. Usar o \textit{Ensaio
sobre a Dádiva}, em \refobra{mauss2017}, como interlocução explícita.}

\atoplanejado{V}{Experimentar o arcabouço semiótico com um caso.
Mostrar como ele descreve mediações e diferenças, quais categorias
emprega e o que permite perceber que a classificação anterior
deixava fora. Tornar o funcionamento de Cecinepá verificável no texto.}

\atoplanejado{VI}{Introduzir uma discordância ou um fracasso do
próprio experimento. Definir como o arcabouço recebe recalcitrância,
assimetria e incomensurabilidade sem pressupor tradução completa
ou consenso de antemão.}

\atoplanejado{VII}{Transformar as condições de participação no bar.
Abrir a passagem para a composição de um mundo comum, sem apresentar
a proposta semiótica como solução suficiente para todas as crises.}

% Seu texto continua aqui. Os titulos dos atos sao opcionais e editaveis.
% \begin{ato}{Titulo do primeiro ato}
% Narrative in English.
% \end{ato}
% Academic prose in English.\fonte[chap.~18]{cunha2017}
